%% =============================================================================
%%  How to do a thematic study (generic: fits any theme)
%% =============================================================================
\chapter{How to do a thematic study}

A thematic study starts from a question (\enquote{\temaPergunta}) and looks
for the answer across the whole Bible. The risk peculiar to this kind of study
is stringing together isolated verses that confirm what one already thought.
The route below exists to avoid that: gather \emph{all} the relevant texts,
read each one in its context, see how the subject develops across the Bible,
and only then frame the synthesis \parencite{osborne2006, duvall2012}.

\section{The seven steps}

\begin{ebtabela}{@{}>{\raggedright}p{.26\linewidth} >{\raggedright\arraybackslash}X@{}}
\EBcab{Step} & \EBcab{What to do}\\\midrule
1. The question & state the question precisely, break it into sub-questions, and record what you already think about the subject (your presuppositions)\\
2. Words and concepts & collect the Hebrew and Greek words of the theme with all their occurrences; look as well for synonyms and for ideas expressed without those words\\
3. The texts & build the study list, group it by subject, and mark the weight of each text (central, support, difficult)\\
4. Reading in context & study the central texts one by one: context, what they say, how they bear on the theme\\
5. Development & see how the theme appears across the Bible: what changes, what stays, and whom each instruction addresses\\
6. Dialogue & face the difficult texts and hear how other translations, commentaries and traditions read them\\
7. Synthesis & answer the question in numbered propositions, each with its textual basis; separate what is clear from what stays open; apply\\
\end{ebtabela}

\begin{alerta}[A word is not a concept]
A theme can be present without the word that names it, and the word can appear
without touching the theme. A concordance is the starting point, not the final
list \parencite{barr1961, silva1994}. In \enquote{Who is Jesus?}, for
instance, the subject lies in titles, actions, relationships and statements,
not in a single word; in \enquote{blood}, the word is in nearly every text,
but not every text treats the same aspect.
\end{alerta}

\begin{alerta}[Cherry-picking]
Choosing only the texts that confirm a conclusion is the commonest mistake in
thematic studies \parencite{carson1996}. That is why the list of occurrences
sits in the appendix: it shows what has not been read yet.
\end{alerta}

\section{How to use this notebook}

\paragraph{The data.} The folder \texttt{\temaDados} holds three text files:
\texttt{textos.tsv} (your study list, which you write and rearrange),
\texttt{vocabulario.tsv} (the words of the theme) and \texttt{ocorrencias.tsv}
(produced by the tool). The commands are in \texttt{MANUAL.md}:

\begin{nota}[Tools]
\small\ttfamily
python3 ferramentas/tema.py ocorrencias H1818 G129 -o \temaDados\\
python3 ferramentas/tema.py vocabulario H1818 G129 -o \temaDados\\
python3 ferramentas/tema.py conferir \temaDados
\end{nota}

\paragraph{Groups and weight.} In the \texttt{grupo} column of
\texttt{textos.tsv} you create your own categories; the notebook groups the
texts by them, in the order they first appear in the file. The \texttt{peso}
column takes \texttt{central}, \texttt{apoio} or \texttt{dificil}; left empty,
the printed notebook shows boxes to decide by hand
({\sffamily\tiny\color{suave}C\,\EBcaixa\;S\,\EBcaixa\;D\,\EBcaixa}).

\paragraph{Text sheets.} To study a text in its context:

\begin{nota}[Example]
\small\ttfamily
\textbackslash begin\{fichatexto\}[group]\{Le 17:11\}\\
\ \ \textbackslash begin\{campo\}\{Context\} ... \textbackslash end\{campo\}\\
\ \ \textbackslash interlinear\{dados/biblia/lv017\_11.hbo.tsv\}\\
\textbackslash end\{fichatexto\}
\end{nota}

\paragraph{Synthesis.} Each conclusion goes into a
\verb|\proposicao{statement}{texts}|. If a text in the basis does not appear
in the study list, that is a sign the list needs revisiting.

\paragraph{Printed or digital.} In \texttt{impresso} mode
(\texttt{config.tex}) the fields you have not filled in yet appear with ruled
lines to write on; in \texttt{digital} mode they disappear.
